% SEGMENT OLP-0428-S01
% Part: normal-modal-logic
% Chapter: axioms-systems
% Section: introduction

\documentclass[../../../include/open-logic-section]{subfiles}

\begin{document}

\olfileid{nml}{axs}{int}

\olsection{அறிமுகம்}

அடிப்படை வாய்ப்புநிலை மொழிக்கு வாய்ப்புநிலை மாதிரிகளைக் கொண்ட ஒரு
பொருண்மையியல் நம்மிடம் உள்ளது. எல்லா மாதிரிகளிலுமுள்ள எல்லா உலகங்களிலும்
மெய்யாக இருக்கும் ஒரு !!a{formula} ஐச் செல்லுபடியாகும் எனும் கருத்தும்,
மாதிரிகள் அல்லது சட்டகங்களின் ஏதோ ஒரு வகுப்புக்குச் சார்பாகச்
செல்லுபடியாகும்---அவ்வகுப்பிலுள்ள எல்லா மாதிரிகளின் எல்லா உலகங்களிலும்,
அல்லது அந்தச் சட்டகத்தை அடிப்படையாகக் கொண்ட எல்லா மாதிரிகளிலும், மெய்யாக
இருக்கும்---எனும் கருத்தும் நம்மிடம் உள்ளன. தருக்கம் பொதுவாக இத்தகைய
செல்லுபடித்தன்மையின் பொருண்மையியல் சிறப்பிப்புகளை, !!{derivability} என்ற
நிறுவல்-கோட்பாட்டுக் கருத்துடன் இணைக்கிறது. ஒரு !!a{formula} செல்லுபடியாக
இருந்தால், இருந்தால் மட்டுமே, அது !!{derivable} ஆகுமாறு ஏதோ ஒரு
முறைமையில் !!{derivability} என்ற கருத்தை வரையறுப்பதே நோக்கம்.

மிக எளியதும் வரலாற்றில் மிகப் பழையதுமான !!{derivation} முறைமைகள்,
ஹில்பர்ட் வகை அல்லது அடிகோள் !!{derivation} முறைமைகள் எனப்படுபவை.
பல வாய்ப்புநிலைத் தருக்கங்களுக்கான ஹில்பர்ட் வகை !!{derivation}
முறைமைகளைக் கட்டமைப்பது ஒப்பளவில் எளிது: மீக்கோட்பாட்டு ஆய்வின்
பொருள்களாகவும் அவை எளியவை (எடுத்துக்காட்டாக, சரித்தன்மையையும்
நிறைவையும் நிறுவுவதற்கு). ஆனால் இயற்கை வருவித்தல் முறைமைகளைவிட இவற்றில்
!!{formula}s ஐ நிறுவுவது மிகவும் கடினம்.

ஹில்பர்ட் வகை !!{derivation} முறைமைகளில், ஒரு !!a{formula} இன்
!!a{derivation} என்பது சில அடிகோள்களிலிருந்து ஒரு சில அனுமான விதிகள் வழியாக
அந்த !!a{formula} ஐ அடையும் !!{formula}s இன் ஒரு தொடராகும். !!{derivation}
முறைமை பொருண்மையியலுடன் பொருந்த வேண்டும் என்பதால், !!{derivable}
வாய்பாடுகளின் கணம் எல்லா மாதிரிகளிலும் மெய்யாகும் (அல்லது எல்லா அடிகோள்களும்
மெய்யாக இருக்கும் எல்லா மாதிரிகளிலும் மெய்யாகும்) என்பதற்கு நாம்
உத்தரவாதமளிக்க வேண்டும். இது செயல்படுவதற்கு அவசியமான வாய்ப்புநிலைத்
தருக்கங்களின் சில பண்புகளை முதலில் தனித்தெடுப்போம்: அவையே “இயல்பான”
வாய்ப்புநிலைத் தருக்கங்கள். இயல்பான வாய்ப்புநிலைத் தருக்கங்களுக்கு மோடஸ்
போனென்ஸ், இன்றியமையாக்கல் ஆகிய இரண்டு அனுமான விதிகளை மட்டுமே எடுத்துக்கொள்ள
வேண்டும். எல்லா மெய்மங்களின் (பதிலீட்டு நிகழ்வுகளையும்), மேலும் நாம் கையாளும்
வாய்ப்புநிலைத் தருக்கத்தைப் பொறுத்து சில வாய்ப்புநிலை அடிகோள்களையும்,
அடிகோள்களாகக் கொள்கிறோம். எல்லா மாதிரிகளின் வகுப்பில் மட்டுமே நமக்கு
ஆர்வம் இருந்தாலும்கூட,~$\Ax{K}$\iftag{notprvDiamond}{}{ மற்றும்~$\Ax{Dual}$}
இன் எல்லாப் பதிலீட்டு நிகழ்வுகளையும் அடிகோள்களாக எண்ண வேண்டும். இவை மட்டும்
குறைந்தபட்ச இயல்பான வாய்ப்புநிலைத் தருக்கம்~$\Log K$ ஐ உருவாக்குகின்றன.

\begin{defn}
\emph{மோடஸ் போனென்ஸ்} விதி பின்வரும் அனுமானத் திட்டவடிவமாகும்:
\begin{prooftree}
\AxiomC{$!A$}
\AxiomC{$!A \lif !B$}
\RightLabel{\MP}
\BinaryInfC{$!B$}
\end{prooftree}
$!C \ident !A \lif !B$ எனில், எனில் மட்டுமே, !!a{formula}~$!B$ ஆனது
!!{formula}s~$!A$, $!C$ ஆகியவற்றிலிருந்து மோடஸ் போனென்ஸால்
\emph{பின்வருகிறது} என்கிறோம்.
\end{defn}

\begin{defn}
\emph{இன்றியமையாக்கல்} விதி பின்வரும் அனுமானத் திட்டவடிவமாகும்:
\begin{prooftree}
\AxiomC{$!A$}
\RightLabel{\Nec}
\UnaryInfC{$\Box !A$}
\end{prooftree}
$!B \ident \Box !A$ எனில், எனில் மட்டுமே, !!a{formula}~$!B$ ஆனது
!!a{formula}~$!A$ இலிருந்து இன்றியமையாக்கலால் பின்வருகிறது என்கிறோம்.
\end{defn}

\begin{defn}
அடிகோள்களின் கணம்~$\Sigma$ இலிருந்து ஒரு \emph{!!{derivation}} என்பது
!!{formula}s~$!B_1$, $!B_2$, \dots,~$!B_n$ இன் ஒரு தொடராகும்; இங்கு
ஒவ்வொரு~$!B_i$ உம் பின்வருவனவற்றில் ஒன்றாகும்:
\begin{enumerate}
\item ஒரு மெய்ம நிகழ்வு; அல்லது
\item $\Sigma$ இலுள்ள ஒரு !!a{formula} இன் பதிலீட்டு நிகழ்வு; அல்லது
\item $j$, $k < i$ ஆக இருக்கும் இரண்டு !!{formula}s~$!B_j$, $!B_k$
  ஆகியவற்றிலிருந்து மோடஸ் போனென்ஸால் பின்வருவது; அல்லது
\item $j < i$ ஆக இருக்கும் !!a{formula}~$!B_j$ இலிருந்து
  இன்றியமையாக்கலால் பின்வருவது.
\end{enumerate}
$!B_n \ident !A$ ஆக இருக்கும் அத்தகைய !!a{derivation} இருந்தால்,
$!A$ ஆனது \emph{$\Sigma$ இலிருந்து !!{derivable}} என்கிறோம்; குறியீட்டில்
$\Sigma \Proves !A$.
\end{defn}

இந்த வரையறையின்படி, !!{derivable} வாய்பாடுகளின் கணம் ஓர் இயல்பான
வாய்ப்புநிலைத் தருக்கமாக அமையும் என்றும், ஒவ்வொரு அடிகோளும் மெய்யாக
இருக்கும் ஒவ்வொரு மாதிரியிலும் எந்த !!{derivable} !!a{formula} உம் மெய்யாக
இருக்கும் என்றும் தெரியவரும். !!{derivation}s இன் இப்பண்பு
\emph{சரித்தன்மை} எனப்படுகிறது. அதன் மறுதலையான \emph{நிறைவு} என்பதை
நிறுவுவது கடினம்.

\end{document}
